\documentclass[twoside]{article}
\setlength{\oddsidemargin}{0.2 in}
\setlength{\evensidemargin}{0.2in}
\setlength{\topmargin}{-0.8 in}
\setlength{\textwidth}{6.5 in}
\setlength{\textheight}{8.5 in}
\setlength{\headsep}{0.75 in}
\setlength{\parindent}{0 in}
\setlength{\parskip}{0.1 in}

%\usepackage{amsmath}
%\usepackage{amsthm}
\usepackage{amssymb}
\usepackage{amsthm}
\usepackage[tbtags]{amsmath}
\newtheorem{thm}{Theorem}
\usepackage[english]{babel}
\usepackage[utf8]{inputenc}
\usepackage{graphicx}
\usepackage[colorinlistoftodos]{todonotes}
\usepackage{hyperref}
\usepackage{setspace}

\begin{document}
\begin{center}
\begin{spacing}{1.75}
{\Large { ECE368: Probabilistic Reasoning} \\
	{Lab 3: Hidden Markov Model}}
\end{spacing}
\vspace{-1em}
\end{center}
%\vspace{-11mm}
You can complete this lab in a group of two. Please provide the name and student number of both members.
\begin{center}
   \framebox{
      \vbox{\vspace{3mm}
    \hbox to 6.28in { {\bf Name:
	\centerline{Student Number:}} }\\[5mm]
    \hbox to 6.28in { {\bf Name:
	\centerline{Student Number:}} }
      \vspace{3mm}}}
   \end{center}
{\bf You should hand in:} 1) A scanned \textsf{.pdf} version of this sheet with your answers (file size
should be under $2$ MB); 2) one Python file \textsf{inference.py} that
contains your code. The files should be uploaded to Quercus.

\begin{enumerate}

\item
\begin{enumerate}
\item  Write down the formulas of the forward-backward algorithm to compute the marginal distribution $p(\mathbf{z}_i|(\hat{x}_0,\hat{y}_0),\ldots,(\hat{x}_{N-1},\hat{y}_{N-1}))$ for $i=0,1,\ldots,N-1$. Your answer should contain the initializations of the forward and backward messages, the recursion relations of the messages, and the computation of the marginal distribution based on the messages. ({\bf $0.5$ pt})
     \begin{center}
   \framebox{
      \vbox{\vspace{40mm}
    \hbox to 5.5in {}
      \vspace{30mm}}
   }
   \end{center}

   \item After you run the forward-backward algorithm on the data in \textsf{test.txt}, write down the obtained marginal distribution of the state at $i=99$ (the last time step), i.e., $p(\mathbf{z}_{99}|(\hat{x}_0,\hat{y}_0),\ldots,(\hat{x}_{99},\hat{y}_{99}))$. Only include states with non-zero probability in your answer.
   ({\bf $1$ pt})

      \begin{center}
   \framebox{
      \vbox{\vspace{25mm}
    \hbox to 5.5in {}
      \vspace{10mm}}
   }
   \end{center}
\end{enumerate}

\item Modify your forward-backward algorithm so that it can handle missing observations. After you run the modified forward-backward algorithm on the data in \textsf{test\_missing.txt}, write down the obtained marginal distribution of the state at $i=30$, i.e., $p(\mathbf{z}_{30}|(\hat{x}_0,\hat{y}_0),\ldots,(\hat{x}_{99},\hat{y}_{99}))$. Only include states with non-zero probability in your answer. ({\bf $0.5$ pt})
    \begin{center}
   \framebox{
      \vbox{\vspace{20mm}
    \hbox to 5.5in {}
      \vspace{10mm}}
   }
   \end{center}

\item
\begin{enumerate}
\item Write down the formulas of the Viterbi algorithm using $\mathbf{z}_i$ and $(\hat{x}_i,\hat{y}_i), i=0,1,\ldots,N-1$. Your answer should contain the initialization of the messages and the recursion of the messages in the Viterbi algorithm. ({\bf $0.5$ pt})
    \begin{center}
   \framebox{
      \vbox{\vspace{40mm}
    \hbox to 5.5in {}
      \vspace{10mm}}
   }
   \end{center}

 \item After you run the Viterbi algorithm on the data in \textsf{test\_missing.txt}, write down the last $10$ hidden states of the most likely sequence (i.e., $i=90, 91, 92,\ldots, 99$) based on the MAP estimate. ({\bf $1.5$ pt})
     \begin{center}
   \framebox{
      \vbox{\vspace{30mm}
    \hbox to 5.5in {}
      \vspace{45mm}}
   }
   \end{center}

\end{enumerate}

\item Compute and compare the error probabilities of $\{\tilde{\mathbf{z}}_i\}$ and $\{\check{\mathbf{z}}_i\}$ using the data in \textsf{test\_missing.txt}.  The error probability of $\{\tilde{\mathbf{z}}_i\}$ is \framebox{\vbox{\vspace{4mm}\hbox to 14mm {}}}. The error probability of $\{\check{\mathbf{z}}_i\}$ is \framebox{\vbox{\vspace{4mm}\hbox to 14mm {}}}. ({\bf $0.5$ pt})
\item Is sequence $\{\check{\mathbf{z}}_i\}$ a valid sequence? If not, please find a small segment $\check{\mathbf{z}}_i,\check{\mathbf{z}}_{i+1}$ that violates the transition model for some time step $i$. You answer should specify the value of $i$ as well as the corresponding states $\check{\mathbf{z}}_i,\check{\mathbf{z}}_{i+1}$. ({\bf $0.5$ pt})
     \begin{center}
   \framebox{
      \vbox{\vspace{20mm}
    \hbox to 5.5in {}
      \vspace{10mm}}
   }
   \end{center}
\end{enumerate}

\end{document}
